\section{后续发展：从 SyntaxNet 到图方法}

\subsection{Google SyntaxNet / Parsey McParseface}

\begin{figure}[H]
\centering
\includegraphics[width=0.95\textwidth]{slides-images/slide-038.jpg}
\caption{Google SyntaxNet 的性能提升\protect\footnotemark}
\end{figure}
\footnotetext{Slides 第38页。}

2016 年，Google 在 Chen \& Manning 的基础上进行了大规模改进，推出了 \textbf{SyntaxNet}（别名 \textbf{Parsey McParseface}）：

\begin{itemize}
    \item \textbf{更深的网络}：使用更多层的神经网络和更大的向量维度
    \item \textbf{更优的超参数}：经过大规模调参
    \item \textbf{柱搜索}（Beam Search）：不再是纯贪心策略，而是同时维护多个候选分析路径，最终选择最优的一条
    \item UAS 从约 92\% 提升到约 \textbf{94.6\%}
\end{itemize}

Manning 提到，Google 将一个依存分析器做成了大新闻——在 Wired、VentureBeat 等科技媒体上广泛报道——这让他感到意外。起了一个搞笑的名字（Parsey McParseface）确实有助于获得媒体关注。

\subsection{神经图方法依存分析（Neural Graph-based Parsing）}

\begin{figure}[H]
\centering
\includegraphics[width=0.95\textwidth]{slides-images/slide-039.jpg}
\caption{图方法依存分析：对每个词评估所有可能的中心词\protect\footnotemark}
\end{figure}
\footnotetext{Slides 第39页。}

除了转移方法，另一种重要的依存分析范式是\textbf{图方法}（Graph-based Parsing）：

\begin{importantbox}{图方法的核心思想}
\begin{enumerate}
    \item 对句子中的\textbf{每对词}，计算一个依存弧的得分（共 $O(n^2)$ 对）
    \item 每个词需要选择一个中心词（或者是 ROOT），评估所有候选中心词的得分
    \item 使用\textbf{最小生成树}（Minimum Spanning Tree）算法从得分矩阵中找出得分最高的合法依存树（无环、连通、唯一根）
\end{enumerate}
\end{importantbox}

\begin{figure}[H]
\centering
\includegraphics[width=0.95\textwidth]{slides-images/slide-040.jpg}
\caption{图方法依存分析的得分矩阵与树构建过程\protect\footnotemark}
\end{figure}
\footnotetext{Slides 第40页。}

图方法的得分计算也可以使用神经网络。Manning 的团队将图方法与神经网络结合，构建了比 Parsey McParseface 更准确的分析器（UAS 再提升约 1\%）。

\begin{figure}[H]
\centering
\includegraphics[width=0.95\textwidth]{slides-images/slide-041.jpg}
\caption{神经图方法分析器的性能对比\protect\footnotemark}
\end{figure}
\footnotetext{Slides 第41页。}

这个神经图方法分析器被集成到了 Stanford 的开源 NLP 工具包 \textbf{Stanza} 中。

\begin{knowledgebox}{转移方法 vs. 图方法}
\begin{itemize}
    \item \textbf{转移方法}：线性时间 $O(n)$，速度快，适合大规模处理；贪心决策可能导致错误传播
    \item \textbf{图方法}：$O(n^2)$ 或 $O(n^3)$ 时间，精度通常更高；考虑全局结构约束（树约束）
    \item 在传统和神经两个时代，图方法的精度都略高于转移方法
    \item 实践中两种方法都被广泛使用，取决于对速度和精度的需求
\end{itemize}
\end{knowledgebox}

\subsection{本章小结}

Google 的 SyntaxNet 通过更深网络和柱搜索在转移方法上取得了进一步提升。图方法通过评估所有词对的依存可能性并求解最优树，提供了另一种精度更高的分析范式。两种方法的神经网络版本都显著优于传统符号化方法。

%% ========== 总结与延伸 ==========

